\begin{abstract}
eBPF lets developers run custom programs inside the Linux kernel, where a static verifier must prove each one safe before it loads. When the proof fails, the verifier rejects the program with a single message that under-determines the repair. The message does not tell the developer which bug occurred, which layer is responsible, or where in the source the fix belongs. On \emph{bpf-bench}, a corpus of 235 reproduced, hand-labeled rejections, one message can stand for many distinct bugs across several responsibility layers, and the flagged instruction is the actual fix site in only 25.5\% of cases. We present \sys, a userspace tool that reads only the verifier log. A rejection is an unfinished safety proof the verifier discards; \sys recovers that proof and localizes the point where a needed safety fact was established and then lost. In a pilot, one model repairs 5 of 18 cases from the raw log and 18 of 18 from the \sys diagnostic.
\end{abstract}
